\section{Discussion}
\paragraph{Information matters more than the classifier}
On 100 patients, selecting among dozens of nearly tied configurations on a single split and reporting the best score adds one to three points of accuracy, about the size of the differences one would use to choose between models or ensemble weights. Under a protocol that removes this bias, changing the classifier, combining several or tuning their weights moves the accuracy by at most a few patients, whereas adding wall-thickness features moves it by 3 to 6 points. With clinically defined groups, the information given to the classifier matters more than the classifier. Several possible sources of bias turned out to be harmless on this data: training on the true LV and testing on a reconstructed one, or selecting features before the cross-validation. They still had to be measured, because each of them could matter on another dataset.

\paragraph{Why the ensembles help only with good features}
On the 21 volumetric features, the MLP ensemble, the SVM and their average all reach about 0.92 to 0.93. Their errors are then largely shared, which leaves little for averaging to correct. With the thickness features, the MLP ensemble and the linear SVM both reach about 0.956, and their average 0.966. The gain of one patient on average is within the noise of our protocol, but it has a consistent sign: on the same splits, the average is better than the linear SVM in 7 repeats out of 10 and than the MLP ensemble in 6, and never worse than either. A linear SVM alone would be an equally defensible choice, and a simpler one to explain.

\paragraph{Comparison with ACDC}
The best ACDC entries reported 92 to 96\% accuracy on the 50 test patients, with features computed on their own automatic segmentations \citep{bernard2018,isensee2017,khened2017}. Our 96.6\% is a cross-validated estimate on the 100 training patients with ground-truth segmentations (and an almost exact reconstruction of the LV), so it is not directly comparable: errors of an automatic segmentation would propagate to the features, and our protocol shares the same 100 patients between model selection and evaluation, even if never within one fold. \citet{isensee2017} also used the whole cardiac cycle, not only ED and ES, which describes the motion of the wall directly.

\paragraph{Limitations}
The test labels are not available to us, so our test predictions cannot be scored, and all our accuracies are estimated on 100 patients. The standard deviations over repeats measure the variability due to the splits, not the variability due to the patients \citep{nadeau2003}; a confidence interval for a new sample of patients would be wider. The choice of the final model among 30 configurations of the benchmark is itself a selection on the same data, which makes its $0.966$ slightly optimistic; the difference with the next configurations is about one patient. The feature sets and the thickness descriptors were designed with the training data in view, without a separate validation set. Finally, the MINF / DCM confusion cannot be resolved by segmentations of two frames: it would require the localization of the lesion or the motion of each wall segment.
